\documentclass[unicode, 12pt, a4paper,oneside,fleqn]{article}


\usepackage{ifxetex}                      %% Для сборки документа и pdflatex'ом, и xelatex'ом
\ifxetex
    %% xelatex
    \usepackage{polyglossia}                       %% загружает пакет многоязыковой вёрстки
    \setdefaultlanguage[spelling=modern]{russian}  %% устанавливает главный язык документа
    \setotherlanguage{english}                     %% объявляет второй язык документа
    \defaultfontfeatures{Ligatures={TeX}}          %% свойства шрифтов по умолчанию
    \setmainfont[Ligatures={TeX}]{Old Standard}    %% задаёт основной шрифт документа
    \setsansfont{Old Standard}                     %% задаёт шрифт без засечек (но Old Standard с засечками заменить на DejaVu Sans?)
    \setmonofont{DejaVu Sans Mono}                  %% задаёт моноширинный шрифт
\else
    %% pdflatex
    \usepackage{cmap}                     %% Поиск русских  слов в pdf
    \usepackage[T2A]{fontenc}             %% Внутренняя кодировка шрифта
    \usepackage[utf8]{inputenc}           %% Кодировка исходного текста;
    %                                      % new: [utf8] (for XeLaTeX);
    %                                      % unmaintained: [ucs];
    %                                      % old, but work: [utf8x] (load ucs)
    %                                      % koi8-r
    %                                      % можно указать cp866 (Alt-кодировка DOS)
    \usepackage[english,russian]{babel}   %% Поддержка русского текста:
    %                                      % включение русификации, русских и
    %                                      % английских стилей и переносов
\fi  



% \hyphenpenalty=50   % переносы разрешены (значение по умолчанию)
% \hyphenpenalty=10000 % запретить переносы
%\hyphenpenalty=9999

% badness - мера разрежённости строки (badness=0 всё хорошо).
% При разбиении абзаца на строки Latex не может создать такие строки,
% badness которых больше, чем значение параметра \tolerance
%\tolerance=10

% к пределу растяжимости каждой из строк в процессе разбиения абзаца
% на строки и вычислений соответствующих значений badness прибавляется
% значение \emergencystretch. Чем больше значение параметра
% \emergencystretch, тем более разреженные строки появятся на печати.
% например: \emergencystretch=10pt конфигурирует TeX для использования
% не более 10pt дополнительного пробела для каждой строки
%\emergencystretch=10pt

% запрет заездов на правое поле документа и переносы. it tells
% TeX to not even look for hyphenation positions
%\pretolerance=9999

% Only words containing at least 16 characters will now be hyphenated.
%\lefthyphenmin=4  % least 4 characters before the hyphen
%\righthyphenmin=4 % least 4 characters after the hyphen


% запрещение вдов ("widow" - короткая строка или слово в конце абзаца)
\widowpenalty=10000
% запрещение сирот ("orphan" - одиночная строка вылетевшая на следующую страницу или колонку)
\clubpenalty=10000
%
% unknown:
%\brokenpenalty=4991
%\predisplaypenalty=10000
%\postdisplaypenalty=1549
%\displaywidowpenalty=1602

\usepackage{graphics}
\usepackage{wrapfig}
\usepackage{multicol}
\usepackage{multirow}
%\usepackage{fullpage}
\usepackage{textcomp} % типографские значки
\usepackage{mathtext} % если нужны русские буквы в формулах
\usepackage{gensymb}  % для спец знаков
\usepackage{amsmath} % для спец знаков в формулах
\usepackage{amssymb} % для спец знаков в формулах
\usepackage{topcapt} % подписи к таблицам
\usepackage{dcolumn} % выравнивание чисел
\usepackage{ulem} % подчёркивание

\usepackage{makeidx} % индекс

\usepackage{import} % for converted svg to pdf

% вращение 
\usepackage{lscape}     % for %\begin{landscape} ...   %\end{landscape}
\usepackage{rotating}   % for sideways and \rotatebox{-90}{}

% водяные знаки
%\usepackage[firstpage]{draftwatermark}
%\SetWatermarkScale{0.7}
%\SetWatermarkLightness{0.9}
%\SetWatermarkText{\textbf{Рабочая версия}}

\usepackage[iso,english]{isodate}
\usepackage{datetime2} % print datetime in GOST ISO 8601-2001

\frenchspacing

\usepackage{fancyvrb} % for verbatim text

\usepackage{listings} % for source code
%\lstloadlanguages{lisp}
\lstset{
  language=C,
  basicstyle=\tiny, %or \small or \footnotesize etc.
  extendedchars=\true, % for russian characters in comments
  %texcl,  % for spaces in russian characters in comments
  keepspaces = true, % for spaces in russian characters in comments. it's work!
  escapechar=|,
  frame=single,
  commentstyle=\itshape,
  inputencoding=utf8,
  stringstyle=\bfseries
}

\usepackage{csvsimple}

% fixme: remove before release
\usepackage[colorinlistoftodos,prependcaption,textsize=tiny]{todonotes}
\setlength{\marginparwidth}{3cm}


\usepackage[colorlinks=true]{hyperref} % url hyperlink (beamer already include it, so move here for prevent conflict)

\author{Роман В. Приходченко}

\title{Руководство по работе с программой}


\makeindex



\begin{document}

% меняем английские термины на русские
\renewcommand\bibname{СПИСОК ЛИТЕРАТУРЫ}
\renewcommand\refname{\centering Список литературы}
\renewcommand\contentsname{\centering Содержание}


% образцы переноса сложных слов - не работает?
% \hyphenation{веб=-ин-тер-фей-се веб-ин-тер-фей-с}
% or use in text: веб"=интерфейс (require: \usepackage[russian]{babel})

% печатаем титульный лист
\makeatletter % generate \@title, \@date, ...
\maketitle

\begin{table}[ht]
  \begin{tabular}{cc}
    \includegraphics[width=2cm]{../../../CC_BY-SA_88x31.png} &
    \shortstack{руководство распространяется в соответствии с
      условиями\\
      \href{http://creativecommons.org/licenses/by-sa/3.0/}{Attribution-ShareAlike} \\
      (Атрибуция — С сохранением условий) CC BY-SA \\
      Копирование и распространение приветствуется.}
  \end{tabular}
\end{table}

\newpage
% печатаем оглавление
\tableofcontents

\section{Алгоритм}
Описание алгоритма и список основных источников в файле\\
\verb'c-artificial-neural-network/doc/growing-neural-gas/ru/algorithm.pdf'


\clearpage
\section{Что содержит внутри себя нейрон и база знаний?}

Юнит-тесты и примеры работы нейронных сетей лежат в каталоге\\
\verb'c-artificial-neural-network/src/example/growing-neural-gas/static-data'
\begin{itemize}
\item \verb'unit-testing.c' --- проверка работы всех функций (можно
  пройтись по всему алгоритму). \todo{Необходима доработка: добавить счётчик
  успешно выполненных функций (сейчас только визуальная проверка).}
\item \verb'gng-to-dot_unit-testing.c' --- проверка работы функций
  вывода в DOT файл (pq.gv)
\item \verb'example.c' --- пример работы маленькой нейронной сети
\end{itemize}

В файле \verb'growing-neural-gas.h' описана структура нейрона,
заметьте, структура содержит ссылки на массивы \verb'weight' и
\verb'conn_age':

\begin{lstlisting}[language=C,
  caption={Описание структуры нейрона},
  captionpos=b,
  label={nueron:code}]
typedef struct Neuron {
	unsigned char active;
        float *weight; // DIMENSION_OF_SENSOR * sizeof (* (gng[i].weight))
	int *conn_age; // LIMIT_NETWORK_SIZE * sizeof (* (gng[i].conn_age))
        float local_error;
	float utility_factor;
	int group; // only for visualization
} NEURON;
\end{lstlisting}


Нейронная сеть это просто массив нейронов (пример из файлов:\\
\verb'unit-testing.c', \verb'gng-to-dot_unit-testing.c',
\verb'example.c', \verb'clusterizator.c').  Здесь выделяется место
для массива нейронов.  Для каждого нейрона дополнительно выделяется
память для массивов \verb'weight' и \verb'conn_age':

\begin{lstlisting}[language=C,
  caption={Выделение памяти},
  captionpos=b,
  label={malloc:code}]
NEURON *pq_gng = malloc(LIMIT_NETWORK_SIZE * sizeof(* pq_gng));
if (pq_gng == NULL) {
	return 1;
}
for (int i=0; i<LIMIT_NETWORK_SIZE; i++) {
	pq_gng[i].weight = malloc(DIMENSION_OF_SENSOR * sizeof (* (pq_gng[i].weight)));
	if (pq_gng[i].weight == NULL) {
		// fixme: free (pq_gng); and already allocation weights and conn_ages
		return 2;
	}
	pq_gng[i].conn_age = malloc(LIMIT_NETWORK_SIZE * sizeof (* (pq_gng[i].conn_age)));
	if (pq_gng[i].conn_age == NULL) {
		// fixme: free (pq_gng); and already allocation weights and conn_ages
		return 2;
	}
}
\end{lstlisting}

И конечно же требуется освободить память в конце работы:
\begin{lstlisting}[language=C,
  caption={Освобождение памяти},
  captionpos=b,
  label={free:code}]
/* free memory */
for (int i=0; i<LIMIT_NETWORK_SIZE; i++) {
	free (pq_gng[i].weight);
	free (pq_gng[i].conn_age);
}
free (pq_gng);
\end{lstlisting}



\clearpage
\section{Как можно прочитать информацию, которую содержит нейрон и база знаний?}

В файле \verb'growing-neural-gas.c' есть функция для отладочной печати
\verb'print_neuron'.  Её можно добавить в файл \verb'clusterizator.c',
но предварительно необходимо уменьшить размер нейронной сети в файле\\
\verb'clusterizator.h' (иначе нейронная сеть просто не влезет в экран)

\begin{lstlisting}[language=C,
  caption={Файл clusterizator.h},
  captionpos=b,
  label={free:code}]
#define LIMIT_NETWORK_SIZE 10
\end{lstlisting}

Также на невлезание в экран --- влияет размер
\verb'DIMENSION_OF_SENSOR', поэтому печать весов (сейчас 1128) можно
отключить в \verb'growing-neural-gas.c' закомментировав
соответствующие строки в функции \verb'print_neuron'.



Следующие добавки могут испортить вывод номера эпохи в
\verb'epoch.txt', и соответственно изменить режим обучения
(решение: использовать\\
\verb'fprintf(stderr, ...)' вместо \verb'printf').

Чтобы добавить в файл \verb'clusterizator.c' отладочную печать:
\begin{enumerate}
\item после считывания из файла базы знаний (содержимое
  knowledge-base.gng) и
\item после расчёта по алгоритму растущего нейронного газа
  (соответственно перед записью в файл \verb'knowledge-base.gng')
\end{enumerate}
В функции main необходимо добавить строки:

\begin{lstlisting}[language=C,
  caption={Файл clusterizator.c},
  captionpos=b,
  label={free:code}]
if ( access(KNOWLEDGE_BASE_GNG_FILENAME, F_OK) != -1 ) { // file exists
	read_gng_from_file (KNOWLEDGE_BASE_GNG_FILENAME, DIMENSION_OF_SENSOR, LIMIT_NETWORK_SIZE, pq_gng);
+	for (int i=0; i<LIMIT_NETWORK_SIZE; i++) {
+		print_neuron (DIMENSION_OF_SENSOR, LIMIT_NETWORK_SIZE, pq_gng[i]);
+        }
	if (epoch == 0L) {
//...........................
extract_groups_from_conn_ages (LIMIT_NETWORK_SIZE, pq_gng);
+for (int i=0; i<LIMIT_NETWORK_SIZE; i++) {
+	print_neuron (DIMENSION_OF_SENSOR, LIMIT_NETWORK_SIZE, pq_gng[i]);
+}
write_gng_to_file (KNOWLEDGE_BASE_GNG_FILENAME, DIMENSION_OF_SENSOR, LIMIT_NETWORK_SIZE, pq_gng);
\end{lstlisting}

\clearpage
\section{Расшифровка отладочной печати}
Результат работы функции \verb'print_neuron' выглядит так:

\begin{lstlisting}[
  caption={Пример отладочной печати},
  captionpos=b,
  label={free:code}]
 веса                        возраст                   локальная   фактор      группа
                            соединений                ошибка      полезности
w:  0.0  -0.6   0.2  -0.4   a: - 0 2 - - - - - - -   e: 0.100   u: 0.600   g: 0
w: -0.0  -1.7   0.8   0.8   a: 0 - 1 - - - - - - -   e: 0.200   u: 0.700   g: 0
w:  0.3  -0.9   0.2  -1.1   a: 2 1 - - - - - - - -   e: 0.300   u: 0.800   g: 0
w:  0.4  -1.2   0.5  -0.7   a: - - - - 3 - - - - -   e: 0.400   u: 0.900   g: 3
w: -0.5   0.2   0.7   0.5   a: - - - 3 - - 0 - - -   e: 0.500   u: 0.000   g: 3
w: -0.3   2.4   4.4   5.3   a: - - - - - - 0 - - -   e: 1.607   u: 7.700   g: 3
w: -0.4   1.3   2.6   2.9   a: - - - - 0 0 - - - -   e: 1.607   u: 8.850   g: 3
-
w:  0.1  -1.2   3.4   5.6   a: - - - - - - - - - -   e: 1.300   u: 2.400   g: 2
-
\end{lstlisting}

Номер группы используются только для картинки и в расчётах не фигурирует.


Таблица возрастов соединений нейронов (нумерация идёт с \O):
\begin{itemize}
\item число обозначает возраст соединения между нейроном в текущей
  строке и нейроном в столбце
\item --- (прочерк) обозначает отсутствие соединения
  (\verb'NOT_CONNECTED') между нейроном в текущей строке и нейроном в
  столбце
\end{itemize}
\begin{lstlisting}[
  caption={Расшифровка таблицы возрастов нейронов},
  captionpos=b,
  label={free:code}]
   0 1 2 3 4 5 6 7 8 9
a: - 0 2 - - - - - - - нейрон №0 соединён с нейронами 1 и 2
a: 0 - 1 - - - - - - - нейрон №1 соединён с нейронами 0 и 2
a: 2 1 - - - - - - - - нейрон №2 соединён с нейронами 0 и 1
a: - - - - 3 - - - - - нейрон №3 соединён с нейроном  4
a: - - - 3 - - 0 - - - нейрон №4 соединён с нейронами 3 и 6
a: - - - - - - 0 - - - нейрон №5 соединён с нейроном  6
a: - - - - 0 0 - - - - нейрон №6 соединён с нейронами 4 и 5
-                      нейрон №7 не существует
a: - - - - - - - - - - нейрон №8 не соединён ни с одним нейроном
-                      нейрон №9 не существует
\end{lstlisting}

\end{document}
